\section{评估模块实现}\label{sec:evaluation_module}

评估模块的目标并不是重复验证上游检测器本身的识别能力，而是围绕 MalBlock 的核心研究问题，对 Agent 在``已有告警证据之后''的决策质量、执行落地情况以及干预后的告警变化进行统一量化。为此，系统在各类实验作业完成后，对作业目录中的结构化产物进行整理与分析，形成统一的评估结果文件。该模块输出的核心结果为 \texttt{evaluation\_report.json}，其中进一步划分为 \texttt{decision\_eval}、\texttt{execution\_eval} 和 \texttt{effect\_eval} 三个部分，分别对应决策有效性、执行有效性和干预收益观察。对于在线回放场景，系统还支持在执行组与无执行组之间进行对照比较，以增强对干预收益的解释力。

\subsection{决策有效性评估}\label{subsec:evaluation_decision}

决策有效性评估用于回答一个基础问题：在已经形成告警窗口和结构化证据后，Agent 输出的高层动作是否总体合理。该部分评估主要面向带标签的实验场景，尤其是 CSV 通道，因为 CSV 通道中的样本标签、攻击类型与统计特征更加稳定，适合用于验证大模型在受限动作空间中的判断能力。

在实现上，系统首先读取聚合后的输入窗口与 \texttt{llm\_decisions.jsonl} 中的决策结果，并依据窗口标识信息完成逐条对齐。随后，评估模块从窗口证据中提取标签、攻击类型、告警统计信息以及决策动作，将每条决策归入统一的分析口径。这里的核心思想不是简单判断``模型说对了没有''，而是将决策拆分为两个层次进行衡量。

第一类是宽口径的风险识别能力，即系统是否能够识别出需要进入干预链路的风险窗口。该部分指标对应 \texttt{risk\_detection\_metrics}，关注的是模型是否对恶意或高风险窗口给出了足够谨慎的动作，而不是立即要求其进入最强处置模式。换言之，这一指标衡量的是 Agent 是否具备将风险样本从普通流量中区分出来的能力，适合反映系统在早期筛查与风险感知层面的表现。

第二类是严格口径的强处置合理性，即在确有必要进行强干预时，系统是否输出了足够强的动作；而在不应强干预时，又是否避免了过度处置。该部分指标对应 \texttt{strong\_mitigation\_metrics}。与前一类指标相比，它更强调 \texttt{block} 等高强度动作的使用是否与样本风险等级相匹配，因此能够反映 Agent 在``识别风险''之外，是否具备``正确选择处置强度''的能力。

这种双层评估设计具有明确意义：如果仅看是否拦截恶意样本，容易掩盖系统过度阻断的问题；如果只看强处置是否准确，又会忽略系统对中间态风险样本的识别价值。因此，系统将风险识别与强处置分开统计，使评估结果更贴近论文中``决策是否合理''这一研究目标。最终，该部分评估不是服务于检测模型比较，而是服务于对 Agent 决策边界、动作分配逻辑与处置策略合理性的分析。

\subsection{执行有效性评估}\label{subsec:evaluation_execution}

执行有效性评估用于回答第二个关键问题：Agent 给出的结构化决策，是否真正被执行联动模块成功转化为系统级处置动作。与决策有效性只关注``该怎么做''不同，执行有效性关注的是``是否真的做成了''，因此它主要面向 Offline 和 Replay 等包含真实执行链路的实验场景。

在系统实现中，每条决策在完成模型推理后，都会经过后处理与工具调用流程，生成包含执行反馈的结果记录。评估模块会读取这些决策记录中的动作、执行模式、\texttt{tool\_result} 以及相关状态字段，并据此统计执行层面的结果。其核心并非只看模型输出了多少条 \texttt{block}，而是进一步区分：哪些动作仅停留在决策层，哪些动作成功触发了执行工具，哪些动作被已有规则覆盖而被系统主动短路，哪些动作虽然输出了高风险决策但最终没有形成新的有效联动。

基于这一过程，系统可得到若干关键指标。其一是工具调用成功数，即联动工具是否返回成功状态；其二是有效执行数，即真正形成了新的处置动作而不是仅在干运行或重复规则下返回成功；其三是重复执行或已覆盖执行数，用于刻画系统在连续窗口中是否能够识别``该目标已经被处置''，从而避免无意义的重复下发；其四是唯一被阻断目标数量，用于衡量实验中实际触发的阻断范围。

这部分评估对于论文尤其重要。原因在于，MalBlock 并不是一个只生成建议的离线分类器，而是一个强调``决策--执行--反馈''闭环的系统。如果缺少执行有效性评估，即使模型输出表面上合理，也无法证明系统具备真实落地能力。通过该部分指标，可以较为清楚地展示：系统是否能够稳定地把高层动作映射为 \texttt{drop}、\texttt{rate\_limit} 或 \texttt{watch} 等具体执行模式，并在有状态执行环境中保持一致、可控和可追踪的联动行为。

\subsection{干预收益观察}\label{subsec:evaluation_effect}

干预收益观察用于回答第三个问题：当系统真正执行阻断、限速或监控动作后，后续告警表现是否发生了有利变化。需要说明的是，这一部分并不试图直接证明``攻击已被根除''，而是以回放流量和告警统计为基础，将干预后的告警压制情况作为一种可观测代理量，从而评估联动处置是否产生了实际收益。

在实现上，系统以每条已执行决策为起点，结合其源地址、执行模式和生效时长，在后续聚合窗口中寻找相同目标的持续告警记录。评估模块会计算处置前后的告警强度变化，并据此形成压制比例等指标。其基本逻辑是：如果某条决策在生效后，相同源地址对应的后续告警显著下降，则说明该次干预具有一定压制效果；如果变化不明显，甚至后续告警继续上升，则说明该次处置的收益有限，或该攻击行为并未被当前策略充分抑制。

为了避免直接使用原始告警数导致结论失真，系统在该部分实现中引入了两项约束。第一，评估以聚合窗口为单位观察前后变化，而不是用全局总告警数直接比较。这样可以使收益观察与具体决策建立对应关系，避免把无关流量的自然波动误当作阻断收益。第二，在 Replay 场景下，系统会对若干典型回放噪声特征进行过滤，例如由高速回放、时间戳异常或 TCP 窗口异常引入的非目标告警，以减少环境噪声对压制比例的污染。对于处置前有效告警数本身过低、无法支撑比值计算的样本，系统会将其剔除，不纳入压制收益统计。

在此基础上，\texttt{effect\_eval} 输出的并不是单一均值，而是包含单条决策级别与整体作业级别两层结果。前者用于观察哪些处置真正产生了明显收益，后者用于汇总整体平均压制水平、可评估样本数量以及异常样本情况，从而支持论文中的案例分析与总体趋势描述。

对于在线回放通道，系统还进一步支持``执行组 / 无执行组''对照评估。其做法是分别读取两组作业的决策结果、回放摘要和评估结果，先检查两次实验在数据源、窗口参数、接口配置和回放环境上是否具有可比性，再对能够匹配上的决策窗口进行成对比较。若两组在关键配置上一致，则进入正式对照；若仅在选中窗口数量等非核心统计项上存在差异，则仅作为警告记录，而不直接判定对照失效。随后，系统比较两组在动作一致性、执行强度、有效执行数、告警压制比例以及过滤后告警量等方面的差异，以判断真实执行是否相对于无执行基线带来了更明显的抑制效果。

因此，干预收益观察的意义不在于构造一个绝对完美的攻击收益定义，而在于为论文提供一条可重复、可解释、与系统执行闭环紧密关联的证据链：即 Agent 的决策不仅能够落地，而且落地之后能够在后续告警层面观察到一定程度的压制变化。

\subsection{模块作用总结}\label{subsec:evaluation__summary}

综合来看，评估模块在 MalBlock 中承担的是``把实验运行结果转化为论文证据''的关键角色。它并不直接参与阻断决策，也不负责上游告警生成，而是在实验结束后对系统输出进行结构化复盘，将原本分散在作业目录中的输入窗口、决策结果、执行反馈和告警变化统一整理为可分析的结果报告。

其中，决策有效性评估回答``决策是否合理''，执行有效性评估回答``动作是否真正落地''，干预收益观察回答``落地后是否产生可观测收益''。三者共同构成了从模型判断到系统联动再到效果反馈的完整证据链。对于本文而言，这一设计能够较好支撑``基于大模型的恶意流量阻断系统''这一主题，使实验分析的重点始终落在阻断决策、执行闭环和干预收益上，而不是偏移为传统意义上的恶意流量检测精度比较。
